\chapter{Introduction}
\label{ch:intro}

\section{Context and Motivation}
\label{sec:intro_context}
In the contemporary landscape of computing, distributed systems, high-performance computing (HPC) clusters, cloud microservices, and enterprise server telemetry infrastructures have reached unprecedented scales of complexity. Modern data centers and distributed software stacks execute millions of concurrent threads across tens of thousands of heterogeneous compute nodes. In such massive environments, runtime component failures, network partitioning, distributed deadlocks, memory leaks, and malicious security breaches are not rare anomalies; they are inevitable operational certainties.

To maintain service level agreements (SLAs), guarantee high availability, and prevent cascading system outages, site reliability engineers (SREs) and operations teams rely heavily on \textbf{system telemetry logs} (syslogs). System logs are ubiquitous, append-only, semi-structured console messages emitted by operating systems, daemons, background services, network interfaces, and user-space applications. Every individual log entry records an instantaneous snapshot of system state, typically comprising a timestamp, hostname or node identifier, severity facility level (e.g., \texttt{INFO}, \texttt{WARNING}, \texttt{ERROR}, \texttt{FATAL}), and an unconstrained natural language textual narrative (e.g., \texttt{"Packet responder 10.0.1.4 timed out after 3000ms"}).

Because logs are generated autonomously without modifying source code or attaching runtime debuggers, they constitute the single most comprehensive and granular diagnostic window into the execution flow of production systems. However, the sheer volume of generated telemetry poses an existential challenge to manual inspection. A modern supercomputer or enterprise cloud cluster routinely produces hundreds of gigabytes of unstructured log text per day—amounting to millions of lines per hour. Human operators cannot manually read or parse this deluge of text to identify root causes. Consequently, there is an urgent and critical demand for automated, highly scalable, and accurate \textbf{log anomaly detection systems} capable of streaming analysis and real-time operational alerting.

\section{The Challenge of Log Anomaly Detection}
\label{sec:intro_challenges}
While anomaly detection is an established sub-field of machine learning, system logs introduce idiosyncratic domain-specific complexities that cause conventional statistical and tabular anomaly detection algorithms to fail:

\begin{enumerate}
    \item \textbf{Unstructured Free-Text Heterogeneity:} Unlike structured numeric metrics (such as CPU utilization percentage or memory consumption), log messages are written by human software developers using arbitrary strings, dynamic variables, IP addresses, hexadecimal memory offsets, and variable parameter values. There is no predefined database schema.
    \item \textbf{High-Volume Streaming and Computational Constraints:} Production monitoring systems operate under severe latency constraints. Telemetry pipelines must ingest, parse, vectorize, and evaluate hundreds of thousands of events per second. Detection pipelines that require computationally heavy deep neural networks, such as multi-layer bidirectional transformers, often incur prohibitively high inference latencies and multi-kilowatt power overheads, rendering edge deployment or localized node monitoring financially and computationally non-viable.
    \item \textbf{Multi-Tenant Concurrency and Asynchronous Interleaving:} In large-scale computing platforms—such as the Blue Gene/L supercomputer with 131,072 processing cores—thousands of independent tasks execute simultaneously. Log messages emitted by distinct nodes, parallel jobs, and system daemons are asynchronously funneled into centralized logging daemons, resulting in an intricately interleaved ``chronological soup'' where sequential dependencies between adjacent lines are severely corrupted.
    \item \textbf{Extreme Class Imbalance and Base-Rate Reality:} Anomalous behavior in well-engineered enterprise systems is rare. In natural operational settings, healthy sequences constitute $90\%$ to $99\%$ of all observations. Machine learning models trained on artificially balanced benchmarks frequently experience the catastrophic ``base-rate fallacy'' when deployed in production, where minor false positive rates unleash an overwhelming flood of false alarms, causing alert fatigue among human operators.
    \item \textbf{Continuous Software Evolution and Semantic Drift:} Software is dynamic. Kernel patches, application updates, and scheduled maintenance periodically introduce previously unseen log templates. Static machine learning models that lack mechanisms to handle out-of-vocabulary events quickly suffer from performance degradation when encountering evolving log distributions.
\end{enumerate}

\section{Research Objectives and Scope}
\label{sec:intro_objectives}
This capstone research project investigates the entire spectrum of log anomaly detection architectures, spanning classical lightweight linear subspace decomposition techniques, robust matrix formulation, continuous natural language embeddings, and non-linear deep neural autoencoders. Rather than treating published literature as unassailable truth, this work adopts a rigorously empirical, reproducible methodology designed to validate, critique, and extend recent findings in the field.

The specific objectives of this capstone research are defined as follows:

\begin{enumerate}
    \item \textbf{Rigorous Baseline Reproduction:} Conduct an exhaustive replication of the linear baseline models proposed in the 2025 IEEE International Conference on System Reliability and Safety (ICSRS) study by F{\"a}lt et al.~\cite{falt2025lightweight} on the structured Hadoop Distributed File System (\textbf{HDFS\_v1}) benchmark. Determine the precise mathematical and preprocessing conditions required to reproduce published Principal Component Analysis (PCA) and Robust Principal Component Analysis (RPCA) results.
    \item \textbf{Algorithmic Discovery of Undocumented Preprocessing:} Identify the hidden failure modes of raw count and TF-IDF representations, and formulate the exact geometric normalization required for linear subspace projection to succeed on structured session data.
    \item \textbf{Mathematical Investigation of the RPCA Reproducibility Gap:} Empirically audit the published $90.55\%$ $F_1$-score claim for Robust PCA via the Inexact Augmented Lagrange Multipliers (ALM) algorithm across multi-seed sweeps, documenting low-rank matrix rank collapse and determining whether the claim is reachable under standard protocols.
    \item \textbf{Exposing the Linear Ceiling on Complex Interleaved Systems:} Apply the reproduced linear frequency models to the massive, multi-node Blue Gene/L (\textbf{BGL}) supercomputer dataset (4.75 million log entries) to quantify the performance degradation of linear models in the presence of massive asynchronous interleaving.
    \item \textbf{Engineering a Deep Semantic Streaming Architecture:} Develop a scalable, high-throughput anomaly detection architecture that integrates physical node-centric temporal windowing, continuous unsupervised Skip-Gram Word2Vec event embeddings, dual-pooling representation (concatenated Mean + Max pooling), non-linear Multi-Layer Perceptron (MLP) autoencoders, and statistical Extreme Value Theory (EVT) Peaks-Over-Threshold (POT) decision boundaries.
    \item \textbf{Execution of a Strictly Leak-Free 10-Fold Protocol:} Design and implement a completely leak-free 10-fold cross-validation protocol reflecting natural operational anomaly prevalence ($10.07\%$), ensuring that every anomalous sequence is tested exactly once, training models exclusively on normal behavior, fitting thresholds strictly without test set observation, and mathematically demystifying earlier project claims of $\sim 90.44\%$ $F_1$.
    \item \textbf{Comprehensive Architectural Ablation and Temporal Robustness Auditing:} Evaluate the necessity of the neural autoencoder against linear PCA and tree-based Isolation Forests on identical feature embeddings, and evaluate chronological temporal degradation under severe Out-of-Vocabulary (OOV) vocabulary shift.
    \item \textbf{Production Efficiency Benchmarking on Consumer Edge Hardware:} Quantify the memory footprint, serialization overhead, parameter count, single-sample latency, batched latency, and end-to-end streaming throughput of the developed pipeline on consumer Apple M2 silicon to assess commercial viability without GPU acceleration.
\end{enumerate}

\section{Structure and Roadmap of the Report}
\label{sec:intro_roadmap}
The remainder of this capstone report is organized into the following chapters:

\begin{itemize}
    \item \textbf{Chapter~\ref{ch:sys_arch}}: \textbf{System Architecture and Methodology (Hardware \& Software Details)} --- Serves as the foundational technical core of the report. It begins with an extensive tutorial background for machine learning practitioners covering the anatomy of syslogs, template induction via Drain, windowing strategies, linear PCA, Robust PCA via Inexact ALM, Word2Vec embeddings, MLP autoencoders, EVT/POT thresholding, and metric evaluation. It then presents the end-to-end streaming system architecture, algorithmic pseudo-code, the empirical reproduction logs of HDFS and BGL, the disclosure of earlier evaluation protocols, and the deployment environment specifications.
    \item \textbf{Chapter~\ref{ch:cost}}: \textbf{Cost Analysis} --- Provides an in-depth computational, hardware, and financial cost audit. It evaluates the parameter efficiency ($5,328$ parameters, $70.25\text{ KB}$), memory cache residency, single-window microsecond latency ($43.2\text{ }\mu\text{s}$), streaming throughput ($62,076\text{ windows/sec}$), and infrastructure OpEx savings of running on commodity CPU hardware compared to enterprise GPU server farms.
    \item \textbf{Chapter~\ref{ch:results}}: \textbf{Results and Discussion} --- Details the complete empirical findings across the 10-fold cross-validation suite under natural prevalence ($10.07\%$). It reports Mode A (label-free POT, $87.80\%$ $F_1$) and Mode B (validation-tuned, $88.66\%$ $F_1$) headline metrics, threshold-free AUC statistics, deconstructs the legacy balanced view ($90.16\%$ $F_1$), presents the architectural ablation study against PCA and Isolation Forest, and analyzes the chronological concept drift failure.
    \item \textbf{Chapter~\ref{ch:conclusion}}: \textbf{Conclusion \& Future Works} --- Synthesizes the experimental insights into the overarching \textit{Dataset-Dependency Rule}, outlines candid operational and engineering limitations, discusses ethical and environmental considerations, and outlines concrete directions for future research.
    \item \textbf{Chapter~\ref{ch:appendix}}: \textbf{Appendix (Source Code)} --- Supplies the core Python source code implementations of the data processing, feature embedding, neural modeling, and POT thresholding engines, along with hyperparameter search grids, complete per-fold confusion matrices, hardware benchmarking scripts, and the project repository file manifest.
    \item \textbf{Chapter~\ref{ch:refs}}: \textbf{References} --- Provides complete, verified bibliographic citations for all foundational literature.
    \item \textbf{Biodata \& End Note} --- Concludes with student researcher profiles and project completion documentation.
\end{itemize}

